
\chapter{Conclusion\label{cha:chapter7}}
\noindent
\indent 
This project focused on the role of serotonin in modulating brain states in anaesthetised mice. Previous evidence suggests that serotonin modulates neural activity, which is assumed to impact the functioning of many cognitive tasks. A specific activity pattern has been observed during anaesthesia and sleep showing periods of high and low activity, also called UP and DOWN states. Experimental investigations showed that transient stimulation of the serotonergic neurons in the DRN induces a state of low activity. To reveal mechanisms behind this phenomenon, a bistable firing rate network was implemented with nodes consisting of a simple E-I circuit, which is simulating the slow-wave oscillations comprising UP and DOWN states. Using data from the Allen Brain Atlas the model is equipped with realistically weighted excitatory long-range connections. The serotonin stimulation is modelled by applying inhibitory currents to the excitatory population. 
\\
\\
\indent Simulations using the proposed model showed that serotonin's effect of inducing a DOWN state could be replicated in all areas, except in the hippocampus and the striatum. Without coupling between the areas, the effect was not observed in areas weakly targeted by the DRN, and increased coupling strength increased the effect of serotonin in those areas. These observations suggest that serotonin acts on an indirect pathway by using inter-area connections to communicate the effect to other areas, depending on the specific connectivity profile of an area. In line with previous studies, this research shows the relevance of long-range connections in communicating neural information. Furthermore, the model simulations demonstrated that modelling the serotonin stimulation as an inhibitory projection targeting the E population leads to the expected effects (DOWN state), whereas the I population with an excitatory current did not result in an induced DOWN state. An UP state was induced when modelling serotonin stimulation was introduced by exciting the E population or inhibiting the I population. that the mechanism of how to model the serotonin stimulation poses an absolute requirement to achieve the replication of the electrophysiological data. 
\\
\\
\indent This project implements an approach that investigates serotonin's effect on a large-scale and therefore provides exciting discoveries about its brain-wide mechanism during anaesthesia. Finding the balance between implementing a biologically plausible model and simplistic network architectures, this project could not answer questions about the mechanisms of serotonin on the level of neurons and receptors.  In future projects, it would be interesting to investigate the role of different receptor and cell types in serotonin’s dynamics. 
\\
\\
\indent In summary, this research emphasises the importance of inter-area communication and highlights the relevance of studies on brain-wide oscillations and serotonin's extensive effects. Moreover, it highlights the great complexity of simplistic neural circuits and unpredictable dynamics when these circuits are interconnected with one another. Understanding the underlying dynamics of serotonin is pivotal as it is the basis of modulating brain activity, which is a key mechanism in many cognitive functions and psychiatric diseases.  

%Idea: make a figure with a sketch of a network with an area that is targeted by "DRN heavy" regions, and therefore has a stronger effect when the projections are selective and it doesn't get influenced by "DRN weak" regions.
